\chapter{Conics}\label{geometry:ch:conics}

\begin{center}\itshape
It is written in the language of mathematics, and its characters are triangles, circles, and other geometric figures.
\upshape(Galileo Galilei, \emph{The Assayer}, 1623, translated by Stillman Drake)
\end{center}

An ellipse and a parabola look different on paper. Seen projectively, each is the same curve with a different line chosen as infinity. That change of viewpoint explains why a single construction can parametrize them, why their tangent lines obey the same rules, and why six points on any of them produce Pascal's line.

Apollonius studied cone sections around 200 BC. Of his eight books, four survive in Greek and three through an Arabic translation; the eighth is lost. Kepler used the word \emph{focus}, or hearth, in his optics of 1604 and found Mars's elliptical orbit in 1609. Galileo's projectile parabola assumes uniform gravity and no air resistance. Pascal's short \emph{Essay pour les coniques}, printed in February 1640 when he was sixteen, recorded the hexagon theorem that we will prove below. These are different reasons for caring about the same family of curves.

We need Appendix~C's projective coordinates and Volume~3's linear algebra. Symmetric matrices will encode conics, but the classification results of the next appendix are not prerequisites: each property used here will be proved here. At the end, two nested circles give a closing problem whose solution comes from choosing the right way to measure arcs.

\begin{think}
Draw a tangent from a point of an outer circle to an inner circle, follow the chord to its other endpoint, and continue consistently in the same direction. If the chain closes from one starting point, must it close from every starting point?
\end{think}

\section{Conics and quadratic forms}

Let $K$ be a field of characteristic different from two, so $2=1+1\ne0$. Examples include $\mathbb Q,\mathbb R,\mathbb C$ and finite fields of odd order.

\begin{definition}\label{geometry:def:conic}
For a nonzero symmetric $3\times3$ matrix $A$ over $K$, put $q(x)=x^{\mathsf T}Ax$ and $\beta(x,y)=x^{\mathsf T}Ay$. Its \emph{conic} $\mathcal C_A$ is the set of projective points $[x]$ satisfying $q(x)=0$. It is \emph{nondegenerate} when $A$ is invertible.
\end{definition}

Homogeneity, $q(\lambda x)=\lambda^2q(x)$, makes the equation independent of representatives. In the chart $z=1$ it becomes $ax^2+bxy+cy^2+dx+ey+f=0$. Degenerate forms can give two lines, as in $xy=0$, or a repeated line, as in $x^2=0$. Some real forms have no projective points at all. The polarization identity
\[
2\beta(x,y)=q(x+y)-q(x)-q(y)
\]
recovers the bilinear form from its quadratic form; this is why characteristic two is excluded here.

\begin{proposition}\label{geometry:prop:conicline}
For a nondegenerate conic $\mathcal C$:
\begin{enumerate}
\item No line is contained in $\mathcal C$, and any line contains at most two points of $\mathcal C$.
\item At $P=[p]\in\mathcal C$, the line $T_P:\beta(p,x)=0$ meets $\mathcal C$ only at $P$. Every other line through $P$ has exactly one further conic point. We call $T_P$ the \emph{tangent}.
\end{enumerate}
\end{proposition}
\begin{proof}
If a two-dimensional subspace $V$ gave a line in $\mathcal C$, then $q$ would vanish on $V$. Polarization would make $\beta$ vanish on $V\times V$, so $V\subseteq V^\perp$. For a basis $v_1,v_2$ of $V$, the two rows $v_i^{\mathsf T}A$ are independent because $A$ is invertible. Their common kernel $V^\perp$ has dimension one, contradicting $\dim V=2$. On any line with basis $u,v$, the binary quadratic $q(su+tv)$ is therefore nonzero and has at most two projective roots. To see the root count including infinity, dehomogenize with $t=1$; if the infinity root occurs, the coefficient of $s^2$ vanishes and the remaining affine polynomial has degree at most one.

Now $Ap\ne0$ makes $T_P$ a line, and $q(p)=0$ puts $P$ on it. On a line through $P$, choose another representative $r$. Then
\[
q(sp+tr)=t\bigl(2s\beta(p,r)+tq(r)\bigr).
\]
If the line is tangent, $\beta(p,r)=0$ and $q(r)\ne0$, otherwise the whole line would be in the conic. Its only root is $t=0$, namely $P$. On a different line $\beta(p,r)\ne0$; the second root $[q(r):-2\beta(p,r)]$ has nonzero second coordinate and gives exactly one more point.
\end{proof}

A quadratic form in three variables has six coefficients up to a common factor. Five points impose five linear conditions. Counting alone suggests uniqueness; the following proof verifies that those conditions really are independent under the stated geometric hypothesis.

\begin{theorem}\label{geometry:thm:fivepoints}
Five distinct points of $\mathbb P^2(K)$ with no three collinear determine exactly one conic equation up to a nonzero scalar. Its conic is nondegenerate.
\end{theorem}
\begin{proof}
Move the first four points to $[1:0:0],[0:1:0],[0:0:1],[1:1:1]$ using a frame map. Write the fifth as $[u:v:w]$. All three entries are nonzero, since a zero entry would put it on a coordinate side. A quadratic through the coordinate points has equation $ayz+bxz+cxy=0$. The last two point conditions are
\[
a+b+c=0,\qquad avw+buw+cuv=0.
\]
The rows are independent: proportionality would give $vw=uw=uv$ up to one common nonzero factor, hence $u=v=w$, repeating the fourth point. Thus the coefficient solution space has dimension one.

Its symmetric matrix has determinant $abc/4$. If $a=0$, the first condition gives $b=-c\ne0$ and the second gives $v=w$. The fifth point would then lie on $y=z$, the line through the first and fourth points. This is forbidden. The same reasoning gives $b,c\ne0$. Hence the matrix is invertible, and moving coordinates back preserves both uniqueness and nondegeneracy.
\end{proof}

\section{Parametrising a conic}

Start with one conic point $P$. Every line through it has a second intersection, except its tangent, which we assign back to $P$. This is a bijection between the pencil at $P$ and the conic. Thus a single known point lets the whole conic be described by a projective-line parameter.

\begin{theorem}[Normal form]\label{geometry:thm:normalform}
Every nondegenerate conic over $K$ that has a $K$-point is projectively equivalent to $y^2=xz$. This conic is parametrized bijectively by
\[
\varphi(t)=[1:t:t^2]\quad(t\in K),\qquad
\varphi(\infty)=[0:0:1].
\]
\end{theorem}
\begin{proof}
Choose a conic point $P$ and a nontangent line through it, giving another conic point $Q$. Such a line exists since a pencil has at least three lines. The two tangents meet at $R$, which cannot be on the conic: the first tangent has only $P$ on it, the second only $Q$. Also $P,R,Q$ are noncollinear. Choose them as the coordinate points in that order. The two zero values $q(e_1)=q(e_3)=0$ and tangent incidences $\beta(e_1,e_2)=\beta(e_3,e_2)=0$ leave
\[
q(x,y,z)=cy^2+2bxz.
\]
Invertibility gives $b,c\ne0$. Substituting $x_{\rm old}=-c x_{\rm new}/(2b)$ gives $c(y^2-x_{\rm new}z)$. Finally a point of $y^2=xz$ with $x\ne0$ is uniquely $[1:t:t^2]$; if $x=0$, then $y=0$ and it is $[0:0:1]$. This proves the normal form and the parameter claim.
\end{proof}

The real circle $x^2+y^2=z^2$ has points and is nondegenerate, so every nonempty nondegenerate real conic is a projective image of it. Over $\mathbb C$ a homogeneous binary quadratic always has a projective root, so every conic has a point and there is just one nondegenerate projective type. Over $\mathbb R$ a nondegenerate form without zeros is definite: two opposite signs on the unit sphere would give a zero along a path between them. Orthogonalizing a basis for this definite bilinear form and rescaling its vectors gives $\pm(x^2+y^2+z^2)$. This proves the remaining empty real type locally; Appendix~E will develop the full classification.

Over $\mathbb Q$, having a point is a genuine restriction. Rational lines through a rational point give all the other rational points. This produces Pythagorean triples from the unit circle, while $x^2+y^2=3$ has no rational point, as the exercises show. The general arithmetic criterion is previewed in Volume~13, not needed here.

\section{Ellipse, parabola, hyperbola}

Choose a real affine chart. Its line at infinity meets a nonempty nondegenerate conic in zero, one or two points; these cases are called an \emph{ellipse}, \emph{parabola} and \emph{hyperbola}. For $x^2/a^2+y^2/b^2=1$, there are no infinity points. For $y=x^2$, homogenization gives $yz=x^2$ and the one infinity point $[0:1:0]$. For $xy=1$, the two infinity points are $[1:0:0],[0:1:0]$, giving the asymptote directions. These are three affine views of the same projective conic.

\begin{proposition}[Focus and directrix]\label{geometry:prop:focus}
Let $F$ be a real point, $D$ a real line not through it, and $e>0$. The locus $PF=e\,d(P,D)$ is a nondegenerate conic: an ellipse for $e<1$, a parabola for $e=1$, and a hyperbola for $e>1$.
\end{proposition}
\begin{proof}
Use Euclidean coordinates with $F=(0,0)$ and $D:x=k$, where $k\ne0$. Both distances are nonnegative, so squaring is equivalent to the original condition and gives
\[
(1-e^2)x^2+y^2+2e^2kx-e^2k^2=0.
\]
Its homogeneous matrix is
\[
\begin{pmatrix}1-e^2&0&e^2k\\0&1&0\\e^2k&0&-e^2k^2\end{pmatrix},
\]
with determinant $-e^2k^2\ne0$. A point exists on the segment from $F$ to its perpendicular foot on $D$, at distance $e|k|/(1+e)$ from $F$. At infinity the equation is $(1-e^2)x^2+y^2=0$, with respectively zero, one or two projective solutions, proving the classification.
\end{proof}

\begin{figure}[ht]
\centering
\fig{geometry-ch05-dandelin}
\caption{Dandelin's spheres, seen from the side. The cutting plane (charcoal) meets the cone in an ellipse; the two spheres inscribed in the cone and
tangent to the plane touch it at the foci $F_1$ and $F_2$.}
\label{geometry:fig:dandelin}
\end{figure}

\begin{beauty}
Dandelin's 1822 construction explains the two foci of an elliptical cone section. Place two spheres inside the cone, tangent to the cutting plane on opposite sides and tangent to the cone along circles. Let $F_1,F_2$ be their plane-contact points. A generator through a section point $P$ meets the sphere-contact circles at $Q_1,Q_2$, with $P$ between them. Tangent segments from the same point to a sphere have equal lengths: the two right triangles formed with the sphere centre have the same hypotenuse and radius. Hence $PF_i=PQ_i$ and
\[
PF_1+PF_2=PQ_1+PQ_2=Q_1Q_2.
\]
The last generator segment has a fixed length between the two contact circles. Thus the section has constant sum of distances to $F_1,F_2$, the focal property of the ellipse. The illustrated elliptical placement is essential to the sum; other cone sections require their corresponding placements.
\end{beauty}

\section{The cross-ratio on a conic, and Pascal's theorem}

Four distinct lines through a point have a cross-ratio: intersect them with any line missing that point and take the cross-ratio of the four intersections. Changing the cutting line is central projection, so Appendix~C proves that the answer is independent of the cut.

\begin{theorem}[Chasles and Steiner]\label{geometry:thm:chaslesconic}
For four distinct points $A,B,C,D$ on a nondegenerate conic, the cross-ratio $(PA,PB;PC,PD)$ is independent of the choice of conic point $P$ different from them.
\end{theorem}
\begin{proof}
Use $\varphi(t)=[1:t:t^2]$. For finite $s\ne t$ the joining line is
\[
\varphi(s)\varphi(t)=[st:-(s+t):1],
\]
by dividing their cross product by $t-s$. Intersect with $x=0$, a line missing $\varphi(s)$, to get $[0:1:s+t]$. Its projective-line coordinate $z/y$ is $s+t$. For $t=\infty$, the joining line meets at $[0:0:1]$, its infinity point. Thus the cross-ratio of the pencil at $\varphi(s)$ is $(s+a,s+b;s+c,s+d)=(a,b;c,d)$, since translation preserves cross-ratios. If $s=\infty$, interchange $x,z$ first. This preserves the conic and replaces parameters by $1/t$, another projectivity, so the same answer follows.
\end{proof}

We can therefore speak of the cross-ratio of four conic points: it is the cross-ratio of their parameters. A related Steiner construction intersects corresponding lines of two projectively related pencils. Choosing two distinct points of this conic as centres and matching the lines with the same parameter recovers the conic, including the tangent at either centre when its parameter occurs. Arbitrary pencil matchings can degenerate, so those hypotheses matter. Coxeter develops the general construction.

\begin{lemma}\label{geometry:lem:perspectivity}
If distinct projective lines $L,M$ meet at $C$ and a projectivity $h:L\to M$ fixes $C$, then $h$ is central projection from a point off both lines.
\end{lemma}
\begin{proof}
Choose distinct $B,U\in L\setminus\{C\}$ and put $V=h(B),D=h(U)$. The joining lines $BV,UD$ are distinct: equality would contain both $B,U$ and hence be $L$, although $V\notin L$. Their intersection $O$ lies off $L$, since otherwise it would equal both $B,U$, and similarly off $M$. Projection from $O$ agrees with $h$ at $C,B,U$, so the frames proposition makes the maps equal everywhere. This is also the criterion proved in Appendix~C's final exercise.
\end{proof}

\begin{theorem}[Pascal]\label{geometry:thm:pascal}
For six distinct points $A,B,C,D,E,F$ on a nondegenerate conic, the three points
\[
X=AB\cap DE,\quad Y=BC\cap EF,\quad Z=CD\cap FA
\]
are collinear.
\end{theorem}
\begin{proof}
At $E$ and $A$, the cross-ratios of the pencils through $B,C,D,F$ agree. Cutting the first pencil by $BC$, and the second by $CD$, gives
\[
(B,C;U,Y)=(V,C;D,Z),\quad
U=ED\cap BC,\ V=AB\cap CD.
\]
Every cut is allowed since a line has at most two conic points. The triples $B,C,U$ and $V,C,D$ are distinct for the same reason. There is a unique projectivity $h:BC\to CD$ taking $B,C,U$ to $V,C,D$. Cross-ratio invariance and injectivity of normalization give $h(Y)=Z$. Since $h(C)=C$, the lemma makes it projection from $BV\cap UD=AB\cap DE=X$. Thus $Y$ and $Z$ lie on one line through $X$, as required.
\end{proof}

\begin{figure}[ht]
\centering
\fig{geometry-ch05-pascal}
\caption{Pascal's hexagram. The hexagon $ABCDEF$ is inscribed in an ellipse; the meeting points $X,Y,Z$ of opposite sides lie on the Pascal line.}
\label{geometry:fig:pascal}
\end{figure}

The tangent at $[p]$ has coefficient vector $Ap$. These tangent lines form the nondegenerate conic $\ell^{\mathsf T}A^{-1}\ell=0$ in the dual plane. Applying Pascal there gives Brianchon's 1806 theorem for six tangents: the three opposite-vertex diagonals are concurrent. The exercise supplies its precise translation. For a pair of base lines, Pascal's incidence pattern becomes Pappus; our nondegenerate proof does not establish that limiting case, which Appendix~C proved directly.

\begin{remark}
The $6!$ orderings, modulo six cyclic shifts and reversal, give $60$ indexed Pascal lines. Special point configurations can make some coincide. For a generic choice the associated triple intersections include twenty Steiner points and further configurations studied by Kirkman, Cayley and Salmon. These incidence results are quoted in Conway and Ryba's account below. Appendix~F will give another proof of Pascal using cubic curves.
\end{remark}

\section{Poles and polars}

\begin{definition}
For a nondegenerate conic $\mathcal C_A$, the \emph{polar} of $P=[p]$ is the line with coefficient vector $Ap$, or equation $\beta(p,x)=0$. Conversely, $P$ is the \emph{pole} of that line.
\end{definition}

Invertibility of $A$ makes polarity a bijection between points and lines. Symmetry gives
\[
P\in\operatorname{polar}(Q)\quad\Longleftrightarrow\quad
\beta(q,p)=0\quad\Longleftrightarrow\quad Q\in\operatorname{polar}(P).
\]
This reciprocal incidence is La Hire's theorem. On the conic itself, the polar is the tangent.

\begin{proposition}\label{geometry:prop:polar}
Let $P\notin\mathcal C$.
\begin{enumerate}
\item A conic point $T$ lies on the polar of $P$ exactly when its tangent passes through $P$. If there are two such tangents, their contact points determine the polar.
\item If a line through $P$ cuts the conic at distinct points $M,N$ and meets the polar at $Q$, then $(M,N;P,Q)=-1$.
\end{enumerate}
\end{proposition}
\begin{proof}
The first assertion is precisely $\beta(t,p)=0$. For the second, $q(p)\ne0$ says $P$ is off its polar, so $Q=[r]$ is a distinct point and $\beta(p,r)=0$. Along this line,
\[
q(sp+tr)=s^2q(p)+t^2q(r).
\]
If $q(r)=0$, its only projective root would be $s=0$, contrary to the two intersections. Hence $q(r)\ne0$. The roots in coordinate $t/s$ are $\lambda,-\lambda$, where $\lambda^2=-q(p)/q(r)\ne0$. In that coordinate $P=0,Q=\infty$, so $(\lambda,-\lambda;0,\infty)=-1$.
\end{proof}

Thus every secant through $P$ supplies a polar point by harmonic conjugation. We have not claimed that such secants exist over every field; an empty real conic has none. For a drawn real ellipse and an exterior point there are plenty of secants. Two of them let a straightedge construct the polar and then its two contact points, as the exercise explains.

\begin{figure}[ht]
\centering
\fig{geometry-ch05-polar}
\caption{The polar of $P$ (rust) passes through the points of contact $T_1,T_2$ of the tangents from $P$. On every line through $P$ that cuts the
conic at $M$ and $N$, the polar cuts at the harmonic conjugate $Q$ of $P$.}
\label{geometry:fig:polar}
\end{figure}

\section{Poncelet's porism}

Let the outer circle $\Gamma$ have centre $o$ and radius $R>0$, and let the inner circle $\omega$ have centre $c$ and radius $r>0$, with $|c-o|+r<R$. At $x\in\Gamma$, choose the directed tangent chord that keeps $\omega$ on its right, and call its other endpoint $T(x)$. Following $T$ means always advancing by the same rule; reversing along the last chord is not allowed.

\begin{theorem}[Poncelet's porism for two circles]\label{geometry:thm:poncelet}
For every positive integer $n$, if $T^n(x)=x$ for one $x\in\Gamma$, then $T^n(y)=y$ for every $y\in\Gamma$.
\end{theorem}

The useful arc density is the reciprocal of the tangent length:
\[
t(x)=\sqrt{|x-c|^2-r^2}>0,\qquad
\mu(I)=\int_I\frac{ds}{t(x)}.
\]
Strict containment keeps $t$ positive on the whole circle, so this gives a finite positive measure to every nonempty arc.

\begin{lemma}\label{geometry:lem:poncelet}
The map $T$ is an orientation-preserving homeomorphism of $\Gamma$, and $\mu(T(I))=\mu(I)$ for every arc $I$.
\end{lemma}
\begin{proof}
We first make the tangent construction smooth. Let $J(a,b)=(-b,a)$, rotation by a positive right angle, and set $v=x-c$, $D=|v|$. The outward unit normal at the chosen contact point is
\[
n=\frac{rv-t(x)Jv}{D^2}.
\]
Indeed $|n|=1$ and $n\cdot v=r$. Put $p=c+rn$ and $d=-Jn$. Direct substitution gives $x=p-t(x)d$. The directed line $x+ud$ therefore touches the inner circle at $p$ and keeps its centre on its right. Since $p$ lies strictly inside $\Gamma$, the other outer-circle intersection is
\[
y=T(x)=x+\ell d,\qquad \ell=-2(x-o)\cdot d>t(x)>0.
\]
All these expressions are smooth on $\Gamma$. Moreover $y=p+t(y)d$. The reverse construction using the other tangent recovers $x$ from $y$, proving bijectivity and hence the homeomorphism claim.

Use counterclockwise arc-length coordinates at both $x,y$. Their positive unit tangents are $u=J(x-o)/R$ and $w=J(y-o)/R$. The chord-endpoint identity gives
\[
n\cdot u=\frac{d\cdot(x-o)}R<0,\qquad
n\cdot w=-n\cdot u>0.
\]
Locally write $n'=\theta'Jn$, differentiating with respect to arc length at $x$. Differentiate $n\cdot(x-c)=r$ and $n\cdot(y-c)=r$, using $x-c=rn+t(x)Jn$ and $y-c=rn-t(y)Jn$. This yields
\[
n\cdot u=-\theta't(x),\qquad
T' n\cdot w=\theta't(y).
\]
The nonzero quantities above allow division, giving
\[
T'=\frac{t(T(x))}{t(x)}>0.
\]
Thus $T$ preserves orientation. Finally change variables on an arc:
\[
\mu(T(I))=\int_I\frac{T'}{t(T(x))}\,ds
=\int_I\frac{ds}{t(x)}=\mu(I).
\]
\end{proof}

\begin{proof}[Proof of Theorem~\ref{geometry:thm:poncelet}]
Choose a basepoint $x_0$ and define a circle coordinate modulo $2\pi$ by
\[
F(x)=\frac{2\pi}{\mu(\Gamma)}\mu([x_0,x]),
\]
where arcs run counterclockwise. Positive continuous density makes $F$ an orientation-preserving homeomorphism. In this coordinate $T$ preserves ordinary arc length. For nearby $a,b$, preservation says that the difference between their image coordinates equals their original difference. Thus $F(T(x))-F(x)$ is locally constant modulo $2\pi$, and connectedness makes it one constant $\alpha$. Therefore $FTF^{-1}$ is rotation by $\alpha$. An $n$th iterate of this rotation fixes a point exactly when $n\alpha\in2\pi\mathbb Z$, and then it fixes every point. Conjugating back proves the porism.
\end{proof}

\begin{figure}[ht]
\centering
\fig{geometry-ch05-poncelet}
\caption{Poncelet triangles. The circles have radii $R=1$ and $r=0.36$, and the distance between their centres satisfies $d^2=R^2-2Rr$. Three
starting points give three triangles, each inscribed in the outer circle and circumscribed about the inner one.}
\label{geometry:fig:poncelet}
\end{figure}

Euler's formula from Appendix~B, $d^2=R^2-2Rr$, is exactly the triangle-closing condition. Its necessity comes from the circumcircle and incircle of a triangle; its sufficiency still requires constructing one triangle, done in the exercises. The quadrilateral condition found by Fuss is
\[
\frac1{(R-d)^2}+\frac1{(R+d)^2}=\frac1{r^2}.
\]
A new exercise proves it, rather than leaving a numerical drawing to support the formula.

\begin{remark}
A broader algebraic Poncelet theorem applies to suitable pairs of nondegenerate complex conics. We quote it here. In the transverse case, the curve of pairs consisting of an outer-conic point and an inner-conic tangent through it is a smooth genus-one curve. Two incidence involutions compose to a translation on that curve; finite order is consequently independent of the initial pair. This is not a single-valued tangent choice on the complex conic itself, nor a claim that every degenerate pair gives a smooth torus. Jacobi's 1828 elliptic-function proof and Cayley's closure criteria are treated by the sources below. Appendix~G introduces cubic models, but the circle proof above uses only real geometry and integration.
\end{remark}

\FloatBarrier
\begin{figure}[ht]
\centering
\fig{v9-appd-pell-parametrization}
\caption{A line of slope $t$ through $P=(0,1)$ meets $y^2=2x^2+1$ again at a rational point when $t$ is rational. The displayed choice $t=1$ gives $(2,3)$. The vertical parameter gives $(0,-1)$; the tangent parameter $t=0$ returns $P$.}
\label{fig:v9-pell-parametrization}
\end{figure}
\begin{figure}[ht]
\centering
\fig{v9-appd-fuss-kite}
\caption{A symmetric Poncelet quadrilateral $ABCD$ with $AC$ a diameter. At $B$ the two tangents are perpendicular. Their unit-normal components give $r^2/(R-d)^2+r^2/(R+d)^2=1$, Fuss's condition. Here $R=1$, $d=3/10$, $r=91/(10\sqrt{218})$.}
\label{fig:v9-fuss-kite}
\end{figure}

Projective coordinates have joined several familiar facts: one conic point supplies a parameter, the bilinear form supplies tangents and polars, and cross-ratio supplies Pascal's line. Poncelet adds a different kind of coordinate, measuring arcs by $ds/t$ so that a moving construction becomes a fixed rotation. The exercises use both ideas before Appendix~E turns to quadratic forms in any dimension.
\FloatBarrier

\section{Exercises}

\begin{exercise}[$\star$]\label{geometry:ex:fivepoints}
Find the real conic through $(0,0),(1,0),(0,1),(1,2),(2,1)$. Classify it and determine its centre.
\end{exercise}
\begin{exercise}[$\star$]\label{geometry:ex:tangentellipse}
For $a,b>0$, find the polar of $(x_0,y_0)$ for $x^2/a^2+y^2/b^2=1$. Deduce its tangent equation when $(x_0,y_0)$ lies on the ellipse.
\end{exercise}
\begin{exercise}[$\star$]\label{geometry:ex:parabolareflect}
Let $f>0$ and consider the parabola with focus $F=(0,f)$ and directrix $y=-f$. For a parabola point $P$, let $H$ be its perpendicular foot on the directrix. Prove that its tangent is the perpendicular bisector of $FH$. Show that a downward vertical ray reflects from it towards $F$.
\end{exercise}
\begin{exercise}[$\star$]\label{geometry:ex:pythagorean}
Use lines through $(-1,0)$ to parametrize every rational point of $x^2+y^2=1$. Deduce the Pythagorean-triple formula $(m^2-n^2,2mn,m^2+n^2)$ up to scaling and the usual changes of sign and order.
\end{exercise}
\begin{exercise}[$\star$]\label{geometry:ex:nothree}
Prove that $x^2+y^2=3$ has no rational solution.
\end{exercise}
\begin{exercise}[$\star$]\label{geometry:ex:pascalpappus}
On distinct lines $\ell,\ell'$, choose triples of distinct points $A,B,C$ and $A',B',C'$, all off their intersection. Verify that Pascal's incidence statement for the hexagon $AB'CA'BC'$ is Pappus. Explain why the nondegenerate Pascal proof alone does not prove this case.
\end{exercise}
\begin{exercise}[$\star$]\label{geometry:ex:brianchon}
State and prove Brianchon's theorem for six distinct tangent lines to a nondegenerate conic: the diagonals joining opposite vertices of the resulting hexagon are concurrent.
\end{exercise}
\begin{exercise}[$\star\star$]\label{geometry:ex:conesection}
A real plane $\Pi:ax+by+cz=1$ cuts $x^2+y^2=z^2$. Prove the section is an ellipse, parabola or hyperbola according as $|c|$ is greater than, equal to or less than $\sqrt{a^2+b^2}$. Interpret the plane as an affine projective chart.
\end{exercise}
\begin{exercise}[$\star\star$]\label{geometry:ex:tangentconstruct}
A real conic is drawn and a point $P$ admits two distinct real tangents to it. Construct both tangents with a straightedge alone, allowing intersections with the drawn curve. Explain your choice of auxiliary secants.
\end{exercise}
\begin{exercise}[$\star$]\label{geometry:ex:concentric}
For concentric circles of radii $0<r<R$, prove that the consistently directed Poncelet map is a rotation. Determine its signed angle with the right-side convention. Show its chain is a convex $n$-gon exactly when $r=R\cos(\pi/n)$, $n\ge3$, and characterize every closing star polygon.
\end{exercise}
\begin{exercise}[$\star\star$]\label{geometry:ex:chappleeuler}
For strictly nested circles of radii $R,r>0$ and centre distance $d$, prove that the Poncelet chain closes after three steps exactly when $d^2=R^2-2Rr$. In that case show every outer-circle point is a vertex of a triangle inscribed in the outer circle and circumscribed about the inner one.
\end{exercise}
\begin{exercise}[$\star$]\label{geometry:ex:apolloniuscircle}
For distinct real points $A,B$ and $k>0$, $k\ne1$, prove that $PA=k\,PB$ describes a circle. Find its centre and radius, and show inversion in it exchanges $A,B$.
\end{exercise}

\begin{exercise}[$\star\star$]\label{ex:v9-finite-conic-tangents}
In $\mathbb P^2(\mathbb F_q)$ with $q$ odd, take $\mathcal C:y^2=xz$. Count its points and tangents. Prove that each off-conic point admits either two or zero tangents, and count the points of each type. For $P=[u:v:w]\notin\mathcal C$, express the two-tangent criterion in terms of whether $v^2-uw$ is a square. Include $u=0$ and the tangent at infinity.
\end{exercise}
\begin{exercise}[$\star\star\star$]\label{ex:v9-pell-conic}
Parametrize every rational point of $y^2=2x^2+1$ by lines through $(0,1)$, including the vertical and tangent parameters. Show that the matrix $A=\left(\begin{smallmatrix}3&2\\4&3\end{smallmatrix}\right)$ preserves this conic and generates all its integer points with $x\ge0,y>0$ from $(0,1)$. Supply a descent proof of completeness, then obtain all other integer points by signs.
\end{exercise}
\begin{exercise}[$\star\star$]\label{ex:v9-poncelet-frequency}
For the strictly nested circles of the porism, put $M=\int_\Gamma ds/t$ and let $\alpha$ be the conjugate rotation angle. Prove that the Poncelet map preserves the probability measure $ds/(Mt)$. If $\alpha/(2\pi)$ is irrational, show every orbit is dense and, for every continuous $h$, its time average equals $M^{-1}\int_\Gamma h\,ds/t$. Deduce the limiting frequency in any arc. Finally bound the normalized arc coordinate's derivative using $R,d,r$ and identify the concentric case.
\end{exercise}
\begin{exercise}[$\star\star\star$]\label{ex:v9-fuss-quadrilateral}
Prove Fuss's quadrilateral criterion for strictly nested circles:
\[
T^4=\mathrm{id}\quad\Longleftrightarrow\quad
\frac1{(R-d)^2}+\frac1{(R+d)^2}=\frac1{r^2}.
\]
Place the centres on the horizontal axis. For necessity, use reflection and the porism to show that an orbit starting at the rightmost point has the leftmost point opposite it. For sufficiency, construct the symmetric kite in Figure~\ref{fig:v9-fuss-kite}. Explain why no orbit has period one or two.
\end{exercise}

\section{Hints}

\begin{hint}{geometry:ex:fivepoints}
Solve for the six quadratic coefficients up to scale, then translate to a critical point of the polynomial.
\end{hint}
\begin{hint}{geometry:ex:tangentellipse}
Multiply $[x_0:y_0:1]$ by the diagonal matrix of the homogeneous equation.
\end{hint}
\begin{hint}{geometry:ex:parabolareflect}
Write $P=(u,u^2/(4f))$. Check a double root on the proposed tangent, then reflect the downward vector rather than just its unoriented line.
\end{hint}
\begin{hint}{geometry:ex:pythagorean}
Substitute $y=t(x+1)$ and factor out the known intersection $x=-1$.
\end{hint}
\begin{hint}{geometry:ex:nothree}
Clear denominators with a primitive integer triple and reduce modulo three.
\end{hint}
\begin{hint}{geometry:ex:pascalpappus}
Write the six sides in order, then pair sides one and four, two and five, three and six.
\end{hint}
\begin{hint}{geometry:ex:brianchon}
Use the conic with matrix $A^{-1}$ in the dual plane and translate Pascal's joins and intersections back.
\end{hint}
\begin{hint}{geometry:ex:conesection}
The infinity directions are cone generators in $ax+by+cz=0$. Rotate the first two coordinates before counting them.
\end{hint}
\begin{hint}{geometry:ex:tangentconstruct}
Four intersections from two secants through $P$ form a quadrangle. Its other two diagonal points determine the polar.
\end{hint}
\begin{hint}{geometry:ex:concentric}
The half-angle $\theta$ of a tangent chord satisfies $\cos\theta=r/R$. Check the direction from the rightmost outer-circle point.
\end{hint}
\begin{hint}{geometry:ex:chappleeuler}
Euler gives necessity. For existence, vary an isosceles triangle continuously from a flat limit to the equilateral triangle, then align its two centres.
\end{hint}
\begin{hint}{geometry:ex:apolloniuscircle}
Expand squared distances and complete the square. Check the two radius vectors have the same direction and product of lengths equal to the squared radius.
\end{hint}
\begin{hint}{ex:v9-finite-conic-tangents}
The finite tangents have coefficients $[s^2:-2s:1]$ and the remaining tangent is $x=0$. Count intersections of pairs of tangents using duality.
\end{hint}
\begin{hint}{ex:v9-pell-conic}
Substitute $y=1+tx$. For integer completeness apply $A^{-1}$ and prove its new nonnegative $x$ is strictly smaller unless $x=0$.
\end{hint}
\begin{hint}{ex:v9-poncelet-frequency}
Transport uniform angular measure through $F$. Use the irrational-rotation result already used in Chapter~2; bounds on $|x-c|$ give bounds on $t$.
\end{hint}
\begin{hint}{ex:v9-fuss-quadrilateral}
Reflection in the centres' axis conjugates $T$ to $T^{-1}$. A four-cycle through a fixed reflection point has its opposite vertex fixed too. The adjacent tangent lines at an endpoint of a diameter are perpendicular at their intersection on the outer circle.
\end{hint}

\section{Solutions}

\begin{solution}{geometry:ex:fivepoints}
\idea Solve the linear interpolation conditions and inspect the translated quadratic.
\writeup Substitution into $ax^2+bxy+cy^2+dx+ey+f$ gives $f=0$, $d=-a$, $e=-c$, $b=-c$ and $a=-b$. Taking $a=1$ yields $x^2-xy+y^2-x-y=0$. Its gradient vanishes at $(1,1)$. With $X=x-1,Y=y-1$ the equation becomes $X^2-XY+Y^2=1$. Since this is $(X-Y/2)^2+3Y^2/4=1$, it is an ellipse centred at $(1,1)$. This also verifies nondegeneracy directly.
\end{solution}
\begin{solution}{geometry:ex:tangentellipse}
\idea Polarity reads the line coefficients from one matrix-vector product.
\writeup The homogeneous matrix is $\operatorname{diag}(a^{-2},b^{-2},-1)$. Multiplying by $(x_0,y_0,1)$ gives the polar $xx_0/a^2+yy_0/b^2=1$. At a conic point the polar is its tangent by Proposition~\ref{geometry:prop:conicline}, giving the same formula for the tangent there.
\end{solution}
\begin{solution}{geometry:ex:parabolareflect}
\idea The equidistance construction identifies the tangent; a vector check determines the reflected ray's direction.
\writeup The focus-directrix equation is $x^2=4fy$. Put $P=(u,u^2/(4f))$ and $H=(u,-f)$. The perpendicular bisector of $FH$ has equation
\[
u(x-u/2)-2fy=0.
\]
It passes through $P$, and substituting into the parabola gives $(x-u)^2=0$, so it is the tangent. Reflection across it exchanges $F,H$. For the direction, its normal is $N=(u,-2f)$. Reflection sends the downward unit vector $v=(0,-1)$ to
\[
v-2\frac{v\cdot N}{|N|^2}N
=\frac{(-4fu,\,4f^2-u^2)}{u^2+4f^2}
=\frac{4f}{u^2+4f^2}(F-P).
\]
The multiplier is positive because $f>0$, so the outgoing ray points towards the focus, including the vertex case $u=0$.
\end{solution}
\begin{solution}{geometry:ex:pythagorean}
\idea A rational slope determines the second intersection rationally.
\writeup Substitution gives
\[
(x+1)((1+t^2)x-(1-t^2))=0.
\]
Thus the second point is $((1-t^2)/(1+t^2),2t/(1+t^2))$. Every rational point other than $(-1,0)$ gives the rational slope $t=y/(x+1)$, and the vertical parameter gives that omitted point. With $t=n/m$, clearing denominators gives $(m^2-n^2,2mn,m^2+n^2)$. Every nonzero Pythagorean triple produces a rational circle point on dividing by its hypotenuse, so this gives the formula up to proportionality, signs and interchanging legs. The exact primitive opposite-parity classification is Volume~2, Theorem~8.1.
\end{solution}
\begin{solution}{geometry:ex:nothree}
\idea A primitive common denominator contradicts divisibility by three.
\writeup A rational solution would give integers $a,b,c$ with $c\ne0$, $a^2+b^2=3c^2$ and $\gcd(a,b,c)=1$. Squares modulo three are zero or one; hence both $a,b$ must be divisible by three. The equation then makes $c^2$ divisible by three and therefore $c$ divisible by three too, contradicting primitivity. No rational point exists.
\end{solution}
\begin{solution}{geometry:ex:pascalpappus}
\idea The alternating vertices lie on a reducible quadratic, but their incidence statement has already been proved directly.
\writeup The six sides are $AB',B'C,CA',A'B,BC',C'A$. Opposite pairs meet at $AB'\cap A'B$, $B'C\cap BC'$, $CA'\cap C'A$, exactly the three Pappus points. Thus the claimed collinearity is Pappus. The quadratic defining $\ell\cup\ell'$ is degenerate, so the conic parametrization and no-three-collinear facts used in our Pascal proof are unavailable. Appendix~C's direct proof establishes this case independently; it is not an automatic application of the nondegenerate theorem.
\end{solution}
\begin{solution}{geometry:ex:brianchon}
\idea The six tangent lines are six points of the dual conic.
\writeup Let the distinct tangent lines be $a,b,c,d,e,f$ in cyclic order. Their adjacent intersections are the vertices. No three of these lines are concurrent, since a line in the dual plane meets its nondegenerate conic in at most two points, so the opposite-vertex joins are defined. Pascal on the dual conic says the dual intersections of $ab$ with $de$, of $bc$ with $ef$, and of $cd$ with $fa$ are collinear. Back in the original plane these three dual points are the lines
\[
(a\cap b)(d\cap e),\quad(b\cap c)(e\cap f),\quad(c\cap d)(f\cap a).
\]
Dual collinearity is their concurrence, exactly Brianchon's conclusion.
\end{solution}
\begin{solution}{geometry:ex:conesection}
\idea An affine section chooses which cone directions count as infinity.
\writeup A point of $\Pi$ represents its line through the origin. Every direction except those in $ax+by+cz=0$ meets $\Pi$ exactly once, so $\Pi$ is an affine chart of the projective plane. The homogeneous circle conic $x^2+y^2=z^2$ appears as the required section in this chart. Rotate about the $z$-axis to replace $(a,b)$ by $(a_0,0)$ with $a_0=\sqrt{a^2+b^2}$. If $c\ne0$, the infinity equation gives $z=-a_0x/c$, so cone directions satisfy
\[
y^2=(a_0^2/c^2-1)x^2.
\]
There are zero, one or two projective solutions as $a_0<|c|$, $a_0=|c|$, $a_0>|c|$. If $c=0$, then $a_0>0$ because the plane equation is nonzero, and the two directions are $x=0,y=\pm z$. These counts give precisely the three classifications. The nondegenerate real projective conic has infinitely many points, so removing at most two infinity points leaves a nonempty section.
\end{solution}

\begin{solution}{geometry:ex:tangentconstruct}
\idea The other diagonal points of an inscribed quadrangle lie on the required polar.
\writeup Choose two distinct secants through $P$ with conic intersections $M_1,N_1$ and $M_2,N_2$. Choose them generically so auxiliary intersections are finite; a drawn ellipse always permits this choice. Put
\[
R=M_1M_2\cap N_1N_2,\quad S=M_1N_2\cap N_1M_2,
\]
and let $C=RS\cap M_1N_1$, $D=RS\cap M_2N_2$. The harmonic construction on the quadrangle $E=M_1,F=M_2,G=N_1,H=N_2$ gives $(R,S;C,D)=-1$. Project from $M_2$ onto $M_1N_1$. The four points go to $M_1,N_1,C,P$, so $(M_1,N_1;C,P)=-1$. Exchanging the last two entries inverts the cross-ratio, still $-1$. Proposition~\ref{geometry:prop:polar} places $C$ on the polar; the same argument puts $D$ there. They are distinct since their secants meet only at $P$, which is off its polar. Thus $RS=CD$ is the polar. Intersect it with the drawn conic at the contact points $T_1,T_2$ and draw $PT_1,PT_2$.

This is a projective construction; if a required point is at infinity its joins are interpreted projectively. For a literal finite drawing the contact points must be finite as well. In the ellipse case this is automatic. In a hyperbola chart an infinity contact means its tangent is an asymptote; the projective construction specifies that line, rather than a finite point to mark on the paper.
\end{solution}
\begin{solution}{geometry:ex:concentric}
\idea Rotational symmetry makes the chord's angular step constant.
\writeup The tangent chord has distance $r$ from the common centre. Its smaller central angle is $2\theta$ with $\theta=\arccos(r/R)\in(0,\pi/2)$. At $(R,0)$, keeping the inner circle on the right chooses the chord heading down and left. Thus $T$ is rotation by $-2\theta$ in the counterclockwise angle convention. A convex $n$-gon is obtained exactly when $\theta=\pi/n$, giving $r=R\cos(\pi/n)$ and $n\ge3$. More generally the chain closes exactly when $\theta/\pi=k/n$ is rational. In lowest terms $0<k<n/2$, its first return is after $n$ steps, winding clockwise $k$ times; it is convex for $k=1$ and a star polygon otherwise. Irrational ratios never return.
\end{solution}
\begin{solution}{geometry:ex:chappleeuler}
\idea Construct one triangle with the correct inradius, then transport it and use the porism.
\writeup A triangle inscribed in $\Gamma$ and circumscribed about $\omega$ has circumradius $R$, inradius $r$ and centre distance $d$. Euler's formula proves necessity. Conversely the equality forces $0<r\le R/2$. In a continuous family of inscribed isosceles triangles, the inradius approaches zero at a flat limit and equals $R/2$ at the equilateral triangle. For example choose vertices $(R,0),R(\cos\theta,\sin\theta),R(\cos\theta,-\sin\theta)$ with $\theta$ from a small positive value to $2\pi/3$; area divided by semiperimeter proves the asserted continuity and limit. The intermediate value theorem supplies inradius $r$. Euler's formula then gives centre distance $d$. A Euclidean isometry aligns the two centres with those of the prescribed circles; if $d=0$, only the common centre needs aligning. Orient the triangle so its interior lies on the right: its tangent steps follow $T$. The porism now closes every orbit in three steps. A step cannot fix a point, because its chord passes through the interior of the outer disk. Nor can two steps return: that would reverse the same chord, placing the inner circle on the left. Hence every such orbit has exactly three vertices and is the required triangle.
\end{solution}
\begin{solution}{geometry:ex:apolloniuscircle}
\idea Completing the square also gives the inversion relation.
\writeup With complex coordinates $a,b,z$, expand $|z-a|^2=k^2|z-b|^2$ and divide by $1-k^2\ne0$. The result is $|z-c|^2=\rho^2$, where
\[
c=\frac{a-k^2b}{1-k^2},\qquad \rho=\frac{k|a-b|}{|1-k^2|}>0.
\]
The vectors $c-a=k^2(a-b)/(1-k^2)$ and $c-b=(a-b)/(1-k^2)$ have the same direction and their lengths have product $\rho^2$. Neither point is the centre. Consequently the inversion formula from Appendix~B exchanges $A,B$.
\end{solution}

\begin{solution}{ex:v9-finite-conic-tangents}
\idea Duality bounds tangent multiplicity; the tangent equation gives the square test.
\writeup The parametrization gives $q+1$ points and distinct tangents. At $\varphi(s)$ for finite $s$, the tangent is $s^2x-2sy+z=0$; at infinity it is $x=0$. Through $P=[u:v:w]$, the finite tangent parameters satisfy
\[
us^2-2vs+w=0.
\]
If $u\ne0$, the nonzero discriminant $4(v^2-uw)$ gives two roots when $v^2-uw$ is a square and no roots otherwise. The discriminant cannot be zero because $P$ is off the conic. There is no infinity tangent in this case. If $u=0$, then $v\ne0$, again because $P$ is off the conic. There is one finite root $w/(2v)$ and the infinity tangent, so two in all; here $v^2-uw=v^2$ is a nonzero square. Scaling the point multiplies this expression by a square, making the test well-defined.

Any two distinct tangents meet off the conic, since a tangent has only its own contact point on the conic. No three are concurrent, by the line-intersection bound for the dual conic. Thus choosing a tangent pair counts each two-tangent point once, giving $\binom{q+1}{2}=q(q+1)/2$. There are $(q^2+q+1)-(q+1)=q^2$ off-conic points, leaving $q(q-1)/2$ with no tangents. This includes every odd $q$, including $q=3$.
\end{solution}
\begin{solution}{ex:v9-pell-conic}
\idea Rational slopes parametrize the curve; a preserving integer matrix admits a decreasing inverse step.
\writeup Substitute $y=1+tx$ to obtain $x((t^2-2)x+2t)=0$. For rational $t$, $t^2\ne2$, so the second point is
\[
x=\frac{2t}{2-t^2},\qquad y=\frac{2+t^2}{2-t^2}.
\]
At $t=0$ it is the double intersection $(0,1)$. The vertical parameter gives $(0,-1)$. Every remaining rational point has $x\ne0$ and the rational slope $t=(y-1)/x$, proving completeness. Figure~\ref{fig:v9-pell-parametrization} shows $t=1$.

The matrix sends $(x,y)$ to $(3x+2y,4x+3y)$; expansion verifies $(4x+3y)^2-2(3x+2y)^2=y^2-2x^2$. Its determinant is one and its inverse sends $(x,y)$ to $(3x-2y,3y-4x)$, preserving integer points too. For an integer solution with $x>0,y>0$, $x=1$ is impossible, so $x\ge2$. The equation gives
\[
\sqrt2\,x<y\le\tfrac32x,
\]
the upper inequality following from $2x^2+1\le9x^2/4$ at $x\ge2$. Hence $0\le3x-2y<x$, and $3y-4x>0$. The inverse step reduces a nonnegative integer $x$ strictly until $x=0$, when $y=1$. Reversing the steps proves every point in this quadrant is $A^n(0,1)$, $n\ge0$. The first are $(0,1),(2,3),(12,17),(70,99)$. Independent sign changes of $x,y$ preserve the equation and produce all other integer points; $y=0$ never occurs.
\end{solution}
\begin{solution}{ex:v9-poncelet-frequency}
\idea The invariant arc coordinate turns the orbit problem into a rotation problem with known uniform averages.
\writeup The lemma gives preservation of $\mu$, so dividing by $M=\mu(\Gamma)>0$ gives an invariant probability measure. The homeomorphism $F$ pushes it to uniform angular measure, since $dF=(2\pi/M)ds/t$. For irrational $\alpha/(2\pi)$ the rotation orbit is dense and equidistributed, by the irrational-rotation result used in Chapter~2 (Volume~5, Corollary~6.14). Apply it to the continuous function $h\circ F^{-1}$ and change variables:
\[
\lim_{N\to\infty}\frac1N\sum_{j=0}^{N-1}h(T^jx)
=\frac1M\int_\Gamma\frac{h}{t}\,ds.
\]
For an arc $I$, approximate its indicator above and below by continuous functions whose integrals differ arbitrarily little. This is possible because $ds/t$ has continuous bounded density and endpoints have measure zero. The resulting visiting frequency is $M^{-1}\int_I ds/t$.

Put $t_- =\sqrt{(R-d)^2-r^2}$ and $t_+=\sqrt{(R+d)^2-r^2}$. Strict containment gives $0<t_-\le t(x)\le t_+$. The normalized coordinate $u=F/(2\pi)$ satisfies
\[
\frac1{Mt_+}\le\frac{du}{ds}\le\frac1{Mt_-},\qquad
\frac{2\pi R}{t_+}\le M\le\frac{2\pi R}{t_-}.
\]
For concentric circles, $t$ is constant, $u$ is ordinary normalized arc length, and the invariant measure is uniform. When $d>0$, $t$ varies, so equal-length arcs need not have equal visiting frequencies.
\end{solution}
\begin{solution}{ex:v9-fuss-quadrilateral}
\idea Reflection forces a diameter in a four-cycle; perpendicular tangent lines then give the metric condition.
\writeup Let $A=(R,0)$, $C=(-R,0)$, and inner centre $I=(d,0)$, choosing any axis if $d=0$. Reflection $\sigma$ in that axis reverses the right-side convention, so $\sigma T\sigma=T^{-1}$. There are no period-one points, since each chosen chord crosses the disk interior, and no period-two points, since reversing a chord puts the inner circle on the wrong side.

If $T^4$ fixes any point, the porism makes $T^4=\mathrm{id}$. The point $T^2(A)$ is fixed by $\sigma$, since $\sigma T^2(A)=T^{-2}(A)=T^2(A)$. It is not $A$ by the period-two exclusion, so it is $C$. Writing $B$ for the upper intermediate vertex, the lower one is its reflection. The quadrilateral is a convex kite, and $\angle ABC=\pi/2$ by the diameter theorem.

Set $a=R-d>r$, $b=R+d>r$. The upper tangent lines from $A,C$ have inward-facing normal directions that, up to an overall sign, may be written
\[
n_A=(r/a,\sqrt{1-r^2/a^2}),\quad
n_C=(-r/b,\sqrt{1-r^2/b^2}).
\]
Perpendicularity of the lines is $n_A\cdot n_C=0$. Both square roots are positive, so squaring this equality is reversible and gives
\[
(1-r^2/a^2)(1-r^2/b^2)=r^4/(a^2b^2)
\quad\Longleftrightarrow\quad r^2/a^2+r^2/b^2=1.
\]
This proves necessity.

Conversely that condition makes the two upper tangents perpendicular. Their intersection is explicitly
\[
B=\left(\frac{2R^2d}{R^2+d^2},\frac{R(R^2-d^2)}{R^2+d^2}\right),
\]
which lies on $\Gamma$ and has positive height because $d<R$. This follows also by solving their equations $bx+ay=bR$, $-ax+by=aR$, since the condition says $r=ab/\sqrt{a^2+b^2}$. Its reflection $D$ gives a convex cyclic kite $ABCD$ containing $I$. The four tangent half-planes $bx\pm ay\le bR$ and $-ax\pm by\le aR$ have precisely this kite as their intersection. Each contains the inner disk, since its boundary has distance $r$ from $I$ and $I$ is on its interior side. Thus every contact point belongs to the disk and the kite, so lies on the actual side segment, not its extension. Traversing it with the interior on the right follows $T$ and closes after four steps. The porism finishes sufficiency. The period exclusions show the first return really has four vertices.
\end{solution}

\section*{Sources and further reading}

Kendig is a friendly continuation of the conic constructions; Coxeter develops their projective setting. For Poncelet, King's paper explains the search for invariant measures, while Dragovi\'c and Radnovi\'c follow the much broader porism and its history. The historical and original records below support the dates and quoted extensions used here.

\begin{itemize}
\item Galileo Galilei, \emph{The Assayer} (1623), in \emph{Discoveries and Opinions of Galileo}, translated by S.~Drake, Doubleday, 1957, pp.~237--238
(the epigraph).
\item Apollonius of Perga, \emph{Conics}: Books I--IV in T.~L.~Heath, \emph{Apollonius of Perga: Treatise on Conic Sections}, Cambridge, 1896;
G.~J.~Toomer, \emph{Apollonius: Conics Books V to VII, the Arabic Translation of the Lost Greek Original}, Springer, 1990.
\item A.~Del Centina, \emph{On Kepler's system of conics in Astronomiae pars optica}, Archive for History of Exact Sciences 70 (2016), 567--589 (Kepler's
foci of 1604).
\item R.~Taton, article ``Pascal, Blaise'' in the \emph{Dictionary of Scientific Biography} (the \emph{Essay pour les coniques} of February 1640 and its
first lemma).
\item G.~P.~Dandelin, \emph{M\'emoire sur quelques propri\'et\'es remarquables de la focale parabolique}, Nouveaux m\'emoires de l'Acad\'emie royale de
Bruxelles 2 (1822), 171--200; K.~Kendig, \emph{Conics}, MAA, 2005.
\item H.~S.~M.~Coxeter, \emph{Projective Geometry}, 2nd ed., Springer, 1987 (Steiner's construction, Pascal, Brianchon, poles and polars);
M.~Berger, \emph{Geometry I and II}, Springer, 1987 (conics, Poncelet).
\item J.~Conway and A.~Ryba, \emph{The Pascal mysticum demystified}, Mathematical Intelligencer 34 (2012), 4--8.
\item J.-V.~Poncelet, \emph{Trait\'e des propri\'et\'es projectives des figures}, Paris, 1822; L.~D.~King, \emph{Three problems in search of a
measure}, American Mathematical Monthly 101 (1994), 609--628 (the invariant measure proof); V.~Dragovi\'c and M.~Radnovi\'c, \emph{Poncelet Porisms and
Beyond}, Birkh\"auser, 2011 (history, Chapple, Euler, Fuss, Jacobi, Cayley); P.~Griffiths and J.~Harris, \emph{On Cayley's explicit solution to
Poncelet's porism}, L'Enseignement Math\'ematique 24 (1978), 31--40.
\item A.~V.~Akopyan and A.~A.~Zaslavsky, \emph{Geometry of Conics}, AMS, 2007.
\end{itemize}
